\subsection{\identifier{UC-08} --- Plan Work}
\label{uc:08}

\paragraph{Purpose}

Allow an authorised user to plan workers,
work areas, optional periods and tasks
for one or more working days.

Planning may start from a new entry,
an existing plan or previously recorded actual work.
It never establishes attendance or actual hours
for the target dates.


\paragraph{Actors}

\textbf{Primary actor:} Farm Owner.

\textbf{Additional actor:} Field Manager,
with the relevant planning permissions granted by the owner.

\paragraph{Preconditions}

\begin{itemize}
    \item The user has permission to plan work
    within the selected scope.
    \item The campaign, work areas
    and worker records exist.
    \item Selected areas are included
    in the campaign coverage.
    \item The requested dates comply
    with the campaign lifecycle rules.
    \item Offline planning requires
    the necessary authorised reference data
    to be available on the device.
\end{itemize}

Planning is optional and supports Planned or In progress campaigns; \identifier{UC-10} does
not require a prior plan. Workers created locally through \identifier{UC-07} may enter pending
plans, with worker acceptance preceding dependent plan acceptance. Planning neither
starts a campaign nor reserves execution.

\paragraph{Main Flow}

\begin{enumerate}
    \item The user selects the campaign
    and one or more target dates.
    \item The user selects a property
    or an included parcel.
    \item The system displays existing plans
    for the selected context.
    \item The user selects the workers
    expected to work.
    \item The user may enter planned periods,
    task types and optional breaks.
    Common values can be applied to a group
    and adjusted individually.
    \item The system checks required information,
    worker availability, area eligibility
    and known scheduling conflicts.
    \item Where a worker is planned
    for several areas on the same day,
    the system requires periods
    for the relevant assignments.
    \item The system presents the proposed
    daily plans and any unresolved issues.
    \item The user reviews and confirms.
    \item The system saves independent daily plans
    and displays their local and server status.
\end{enumerate}

\paragraph{Alternative Flows and Exceptions}

\begin{description}

    \item[A1 --- A minimal plan is sufficient.]
    For a worker assigned to one area
    on a working day, the user may save
    the campaign, date, area and worker
    without specifying hours or task type.

    \item[A2 --- Individual details differ.]
    The user adjusts one worker's periods
    or tasks without unintentionally changing
    other workers' entries.

    \item[A3 --- Several areas apply.]
    The user creates the relevant assignments.
    Periods are required for all assignments
    of that worker to different areas
    on the same working day.

    \item[A4 --- Several target dates are selected.]
    The system applies the reviewed proposal
    to each selected date.
    Each date receives an independent plan,
    validated against its own context.

    \item[A5 --- An existing plan is copied.]
    The user selects a source plan
    and the target dates.
    Workers, areas, periods and tasks
    are proposed for review.
    Source plans remain unchanged.

    \item[A6 --- Actual work is used as a basis.]
    The user selects previously recorded work.
    Its workers, areas, periods and tasks
    are proposed as planning information
    for the target dates.
    Confirmed attendance and actual work records
    are not copied to those dates.

    \item[A7 --- Target dates already contain plans.]
    The system identifies the existing plans.
    Reuse does not silently overwrite
    or duplicate them.
    Changes require explicit review.

    \item[A8 --- A worker is inactive.]
    The system identifies the inactive worker
    and requires an eligible selection
    or authorised reactivation.
    Existing future plans are flagged
    rather than silently edited.

    \item[A9 --- Periods overlap.]
    The system blocks conflicting periods,
    considering the worker's relevant plans
    across properties, parcels and campaigns.

    \item[A10 --- Work crosses midnight.]
    The user records the start and end
    with sufficient date information.
    The period belongs to the working date
    on which it starts.
    Validation considers the full interval,
    including its continuation into the next day.

    \item[A11 --- A worker is removed or replaced.]
    Before actual work has been recorded
    for the affected daily area context,
    the user may confirm removal or replacement.
    Other assignments remain unchanged.

    \item[A12 --- Actual work already exists.]
    The corresponding plan remains
    the reference for what was expected.
    Differences are recorded as actual work
    or handled through the work-correction process.

    \item[A13 --- The device is offline.]
    The user may create or adjust plans
    using prepared authorised data.
    The system identifies the changes
    as pending synchronisation.

    \item[A14 --- Planned breaks are entered.]
    The user identifies break start and end times
    and their intended counting treatment.
    The system prevents duplicate deductions
    in any displayed planned-duration calculation.
    These values remain subject to actual confirmation.
\end{description}

\paragraph{Postconditions}

\textbf{On success:}
Daily plans identify the selected campaign,
working dates, areas and workers,
with any supplied periods and tasks.

\textbf{On reuse:}
Target plans are independent of their source.
Editing a target day does not modify
the source or other target days.

\textbf{On pending synchronisation:}
Local plans remain distinguishable
from server-accepted plans.

\textbf{On cancellation or rejection:}
Previously accepted information remains unchanged.

\textbf{In all cases:}
Planning does not confirm attendance,
create actual hours or modify source actual work.

Period validation follows \identifier{FR-026}; counted-time interpretation is defined in Section~\ref{sec:counted-working-time}.

\paragraph{Related Requirements and Rules}
\hyperref[fr:024]{\identifier{FR-024}}, \hyperref[fr:025]{\identifier{FR-025}}, \hyperref[fr:026]{\identifier{FR-026}}, \hyperref[fr:027]{\identifier{FR-027}}, \hyperref[fr:029]{\identifier{FR-029}}.

\hyperref[br:015]{\identifier{BR-015}}, \hyperref[br:016]{\identifier{BR-016}}, \hyperref[br:017]{\identifier{BR-017}}, \hyperref[br:019]{\identifier{BR-019}}, \hyperref[br:021]{\identifier{BR-021}}.

\paragraph{Decision References}
\decisionref{Q-WF-04}, \decisionref{Q-AA-04}, \decisionref{Q-WF-02}.
